% 迁移来源：NSFC_2026_CIESM-STSR/extraTex/2.1.研究内容.tex
\subsubsection{研究任务与总体安排}

四项研究依次建立训练依据、联合模型、受约束输出与应用证据，共用数据划分和评价标准。参考资料内开展配对训练，真实 CIESM 用于非配对应用，测试集不参与调参（图~\ref{fig:framework}）。
\CiesmFigure{figures/F02-framework.pdf}{四项研究及训练—应用关系：参考资料内配对训练，真实 CIESM 独立应用，测试集不参与方法反馈}{fig:framework}

\paragraph{研究内容一：CIESM 多源多尺度气象数据构建与目标场信息含量检验}
核对 CIESM、ERA5、ERA5-Land、CMFD、台站日值及地形资料的变量定义、单位、日历、时间聚合和空间支撑。\textbf{逐变量检验目标场相对上游强迫的信息增量}，以仅插值对照筛选参考源，初步采用 ERA5-Land 气温及 CMFD 风、短波和降水。选源后固定粗细训练对、CIESM 非配对样本、独立台站验证集和地形分层，形成来源、覆盖与信息增量记录。

\paragraph{研究内容二：多变量多尺度时空超分辨率方法}
采用共享时空表征与变量专用输出，比较联合模型和等信息单目标模型，辨别四变量共同变化与专有变化。以日尺度空间重建和邻日连续性为核心，检验时间上下文相对单日输入的增益；无法唯一确定的细节由可选条件分布描述。

\paragraph{研究内容三：复杂地形引导与尺度一致性约束}
分析高程差、坡度坡向和下垫面的作用，分别处理聚合、非负及变量定义关系与经验性梯度、依赖约束。比较精度、稳定性和极端保持，识别过度约束及域外失效，形成可逐项诊断的 TerraSR-XJ 投影模块。

\paragraph{研究内容四：CIESM 集成、台站独立验证与区域应用}
集成区域增强模块，检查公开资料兼容性，取得实验室接口后联调。在完整地形单元和连续独立时段评价均值、时空结构、联合分布与极端；以\textbf{不参与训练的国家级台站日值}独立验证，并与原始 CIESM、校正场和插值基线比较。验收依据为实验、接口运行及实际应用记录。

\subsubsection{研究目标}

形成可解释粗细尺度关系的 CIESM 区域增强方法，明确重建收益、约束作用与适用边界，围绕 P02“算法或数据集”交付两项成果。

\paragraph{成果一：TerraSR-XJ 算法模块}
开发可独立于网络复用的开源模块：输入粗网格场、面积算子和地形条件，经可行性检查、约束投影及逐变量非退化门控，输出满足聚合、非负和变量定义关系的细网格场及逐网格残差、门控记录。相较规整像素块均值修正\cite{Geiss2022}，重点处理非嵌套网格、边界部分覆盖和算子秩亏；不可行时按预定规则降级并记录失败。随附接口文档、运行样例与单元测试。

\paragraph{成果二：新疆区域高分辨率试验数据产品}
范围为 73$^{\circ}$E—96$^{\circ}$E、34$^{\circ}$N—50$^{\circ}$N 陆地掩膜区，采用 0.1$^{\circ}$ 逐日网格，包含近地面气温、近地面风、地表向下短波辐射和降水。覆盖期取 CIESM \texttt{historical} 与参考资料交集，以测试期为交付区段；采用符合 CF 约定的 NetCDF，附来源、频率、粗化算子、误差、台站验证和使用限制。产品标注为模型推断结果，按上游许可与实验室约定发布。

统一记录四变量来源与处理过程，完成文件读入、细网格输出、粗化复核和应用示例。唯一主指标为配对实验的无量纲地形分层空间结构误差，关键次指标为跨变量秩相关矩阵误差，并报告逐变量聚合残差、非法输出比例及极端表现。指标权重、对照和阈值在训练、验证阶段固定；MCID 与非退化规则沿用立项依据，不将预期改善写成既得或必达效果。

\subsubsection{拟解决的关键科学问题}

\paragraph{其一，粗细尺度关系在何种条件下可迁移？}
参考资料与 CIESM 的偏差和天气轨迹不同，空间粗化、季节与输入分布如何影响重建，校正后哪些气候信息仍可保留和验证？

\paragraph{其二，地形条件如何改变变量共享的收益？}
共享可能利用变量相关性，也可能受高方差变量主导；相对等信息独立基线，其收益如何随地形、季节和极端样本变化？

\paragraph{其三，一致性约束怎样影响有效细节？}
聚合匹配不保证细节正确，经验约束也可能抑制真实异常；如何区分数值一致性与独立样本收益，并确定失效时的简化方案？

\subsubsection{研究方法与技术路线}

以下四个环节依次对应研究内容一至四：固定资料、划分和聚合算子，建立联合模型，在同一模型与样本上分析地形及约束，最后以选定配置开展 CIESM 验证与集成。各环节共享版本化数据，测试集不参与方法修正。“时空”限于日尺度空间重建及邻日连续性；日内（6 小时及更高频）的山谷风、辐射循环重建需要高频参考和更高算力，留作后续研究，一年期资源集中于四变量日尺度重建、等信息对照和集成验证。

\paragraph{数据构建与参考场检验}
核验公开 CIESM \texttt{historical} 与参考资料共同期，记录版本、试验、成员、网格、频率、校验值及变量统计定义。风优先取东西、南北分量，仅有风速时不评价风向和向量关系。统一气温单位，按时间边界积分降水率、按累计区间换算辐射，核对闰日、日历、累计起点和时区。

训练、验证、测试期分别为 1981—2005、2006—2009、2010—2014 年；覆盖不足时按同比例前移并记录。先划分时空集合，再估计标准化、校正、地形统计和阈值；时间边界留出至少半个窗口，空间按完整地形区分组，训练增强不进入留出区。台站验证同期于测试集，站点不参与训练和调参。

以反映 CIESM 面积和采样方式的算子 $A$ 将参考场 $Y$ 粗化，构造 $(AY,Y)$。气温初选含高程订正的 ERA5-Land，其余三变量初选融合台站、卫星与再分析信息的 CMFD\cite{He2020CMFD}，各源均须检验信息含量：在训练集拟合轻量等信息独立模型，以验证集“仅插值 MAE 减去模型 MAE”为增益，用模型验证集 RMSE 按立项依据确定逐变量 MCID。95\% 区间下界高于阈值才通过，上界低于为未达有意义改善，其余为证据不足；后两类均不采纳，选源在联合建模前固定，不接触测试集。

ERA5 同日粗场仅在 CIESM 有同定义、同频率对应项时作为额外条件，记录其与目标资料的关系，不引入应用时不可得信息。真实 CIESM 样本 $X$ 不与 $Y$ 逐日配对。面积权重按陆地掩膜及粗细网格共同覆盖区确定，仅对完整覆盖、权重成立的粗格严格聚合，边界和缺测另报。样本索引、预处理程序及差异说明经手工抽查、粗化复算和隔离检查后，固定供后续环节使用（图~\ref{fig:data}）。
\CiesmFigure{figures/F03-data.pdf}{资料组织与样本隔离：两类资料共用训练、验证、测试划分约定，来源与集合不一一对应}{fig:data}

\paragraph{分布校正与时空联合重建}
以训练期 $AY$ 分布拟合变换 $B$，使 $X^{\star}=B(X)$ 与训练条件的变量定义、尺度支撑一致。比较不校正、逐变量分位数及趋势保持校正，以多变量统计校正作强化对照，按验证集边缘分布和依赖表现选取\cite{Lange2019ISIMIP,Cannon2018MBCn}，不以测试目标重拟合。跨时期分别报告 $X$、$X^{\star}$ 与重建场变化，不将校正信号视为原模式信号原样保留。建模前比较 $AY$ 与 $B(X)$ 的空间功率谱、自相关尺度和邻日自相关；超过验证资料预定容差时作输入侧谱域匹配或随机粗化增强，单报其影响，不归为模型自身收益。

模型写为
\begin{equation}
\widetilde{Y}_{t}=G_{\theta}\bigl(X^{\star}_{t-k:t+k},Z,S_t,\xi\bigr),
\end{equation}
其中 $Z$ 为地形和下垫面，$S_t$ 为季节或日序，$\xi$ 为可选随机变量。用多尺度卷积或轻量时序模块刻画背景与局地残差，设置变量专用输出，按训练尺度归一损失；先学条件中心场，再检验小规模随机残差对概率校准和空间分布的增益。气温、风的连续误差与降水发生、正值量分别建模，独立评价零值频率和极端尾部。连续 7 日输入与单日、邻日乱序和短窗口对照；前后邻日仅用于离线重建，不代表实时预测。

对照由面积一致插值、统计校正加插值，递进至独立模型、共享模型、联合加地形及显式约束模型。独立模型使用相同多变量粗场和地形，总参数量级、预算一致；传统单变量输入仅作应用参照。共享/独立与有/无地形构成两因素试验，固定校正和投影规则；交互成立才解释为地形增强共享收益，仅共享主效应成立则限于共享收益。依据验证集主次指标与时间对照选定中心场、可选集合及配置，供约束分析使用（图~\ref{fig:reconstruction}）。
\CiesmFigure{figures/F04-reconstruction.pdf}{共享表征与四类专用输出；随机残差为可选扩展，独立基线保持同信息和总预算}{fig:reconstruction}

\paragraph{地形引导与尺度一致性约束}
TerraSR-XJ 以候选场、聚合算子和地形分层为输入。共同有效覆盖区内粗格 $c$ 内细格 $i$ 的面积权重 $w_{ci}$ 满足 $\sum_iw_{ci}=1$；对允许有符号残差的连续变量，研究
\begin{equation}
\widehat{Y}=\underset{y}{\arg\min}\;\|y-\widetilde{Y}\|_{W}^{2},
\qquad \text{满足}\quad Ay=X^{\star},
\end{equation}
其中 $W$ 为正定权重矩阵，$A$ 固定。非嵌套网格先查可行性、秩和数值条件，不套用规整像素块均值修正，并输出逐格投影残差。

嵌套网格降水按非负细节权重 $q_i$ 归一：
\begin{equation}
\widehat{P}_i=P_c^{\star}\frac{q_i}{\sum_{j\in c}w_{cj}q_j}.
\end{equation}
要求 $P_c^{\star}\geq0$ 且分母为正；粗格零值则细格全零，分母退化则非负均匀分配并记录，不以隐含常数扰动总量。发生掩膜与正值权重分开，避免毛毛雨扩散；非嵌套网格使用非负数值投影并查可行性。短波同样检查非负和相应时段聚合。

日均风分量模长不能替代日平均风速；只有同瞬时分辨率分量可先合成再日聚合，仅有日均标量风速时独立建模。上述条件仅检验变量定义和聚合关系，缺少蒸发、储水、辐射净收支及全部热通量时，不代表完整水量、能量或动量守恒。

地形条件包括模式与细网格高程差、起伏、坡度坡向及陆面类别；温度—高程关系按季节和地形估计，允许逆温。梯度和变量依赖采用软惩罚，权重仅由验证集选定；从零递增，结合等预算消融、干湿分层及冲突诊断，比较主次指标、极端、物理范围异常率和一致性残差，报告收益—代价曲线及失效条件。数值一致性与独立样本收益分别检验，选定配置进入 CIESM 应用（图~\ref{fig:constraints}）。
\CiesmFigure{figures/F05-constraints.pdf}{聚合、非负和变量定义条件分别检查；经验约束另行分析，数值一致性与独立样本收益分开评价}{fig:constraints}

\paragraph{多层评价与 CIESM 集成验证}
配对实验报告四变量面积加权 MAE、RMSE 及异常相关，误差列物理单位与训练尺度归一值。主指标取地形分层半变异函数无量纲误差：距离档绝对差除以训练期参考半变异幅度，再按距离档、变量、地形加权（非负、各层和为 1）；近常数层设正尺度下限并单报。次指标取去季节、训练尺度标准化的变量秩相关矩阵误差，同测试集对照等信息独立基线。MCID、零基线与非退化规则见立项依据；逐变量 MAE 用物理单位 MCID，区间下界高于/上界低于阈值分别为支持/不支持，其余证据不足。尺度、档位、权重和阈值在训练、验证期确定，测试前冻结。

补报功率谱、邻日自相关、湿日频率、强降水分位及高温/大风持续性。四变量复合事件按训练期分位数定阈值，高值取第 95 百分位、湿日取 1\,mm/d，检验测试期共现及阈值敏感性；持续天数不拼接缺测。按独立数字高程模型（DEM）和地理范围整区留出天山及南北坡、准噶尔与塔里木盆地，避免小块重叠泄漏；结论限区域内留出及台站验证。

验证站不参与训练、调参或 CMFD 生成；同站点集按网格取值并按空间代表性加权，对比测试期国家级台站与原始、校正、重建及插值场的气候态、季节循环、分位数、极端频次，不逐日强配对。台站、参考资料配对结果分报，ERA5/ERA5-Land 互比不作独立观测证据。真实 CIESM 另比原始、校正、增强场的季节均值、分位数、空间谱、联合统计及向 $X^{\star}$ 回聚合残差；相对插值的目标一致性提高不等于精度优于原模式，须凭台站独立验证。

随机模型报连续排序概率评分、区间覆盖与宽度，对照训练期气候分布和确定性模型残差集合。训练 3 个独立种子，按年份、事件时空块重抽样求 95\% 区间，种子及相关格点不作独立样本；年份较少时逐年列示，校正多重比较。主指标检验假说，次要及分层指标作解释，不将局部改善外推全区。

固定模型完成 CIESM 读入、分块推断、细网格输出、粗化复核及接口联调，交付模块、数据和示例，记录环境、资源、变量语义；各阶段证据要求见图~\ref{fig:integration}。
\CiesmFigure{figures/F06-integration.pdf}{拟开展的 CIESM 集成流程；公开兼容性、联调、实际应用与接收均须对应的真实记录}{fig:integration}

\subsubsection{可行性分析}

公开模式与再分析资料、统计降尺度、卷积重建及投影方法提供可复现起点；申请人的多尺度时间序列研究积累和现有 GPU、软件环境支持区域分块训练，详见“研究基础与工作条件”。

第一季度核实四变量共同频率、目标场信息量及实验室接口。某变量所有候选源均未通过信息检验时退出重建，仅保留分布校正和一致性检查，并如实报告未支持有意义重建；关键变量只有月资料时形成月尺度阶段成果，日尺度核心验证仍记为未完成；生成模型超出预算时优先完成确定性模型、约束消融和接口验证，概率扩展留待后续。四种比较配置按各类资料的评价口径检查非退化与不确定性（图~\ref{fig:evaluation}）。
\CiesmFigure{figures/F07-evaluation.pdf}{四种拟比较配置与两类验证证据；分别评价，并检查非退化条件和不确定性}{fig:evaluation}
